% =============================================================================
\chapter{Podsumowanie i dalsze prace}
\label{ch:conclusions}
% =============================================================================

Zapytaliśmy, czy paradygmat programowania wpływa na kod pisany przez LLM. Aby odpowiedzieć na to
pytanie, zbudowaliśmy benchmark dziesięciu zadań algorytmicznych z identycznymi ukrytymi testami
w siedmiu językach z czterech paradygmatów, wraz z systemem, który automatycznie generuje, testuje,
naprawia i ocenia rozwiązania. Pełny przebieg modelu \haiku{} daje następujący obraz:

\begin{itemize}
  \item \textbf{Poprawność zależy bardziej od języka niż od paradygmatu.} 80\,\% zadań rozwiązano za
    pierwszym podejściem, a 93\,\% po samonaprawie. Różnice między dwoma językami jednego
    paradygmatu są równie duże jak różnice między paradygmatami i żadna nie jest istotna przy
    jednej próbce na zadanie. Prolog wypada najsłabiej, częściowo z powodu dwóch defektów w jego
    zestawach testów.
  \item \textbf{Samonaprawa jest skuteczna, ale ograniczona przez informację zwrotną.} Dwie rundy
    naprawy usunęły 64\,\% błędów, a trzecia runda nic nie dodała. Naprawa utykała tam, gdzie
    informacja zwrotna nie ujawniała przyczyny.
  \item \textbf{Idiomatyczność zależy od paradygmatu.} Kod funkcyjny był najbardziej idiomatyczny,
    a obiektowy i logiczny najmniej. Oba języki każdego paradygmatu są tu zgodne, więc jest to
    najbardziej spójny efekt paradygmatu, jaki zaobserwowaliśmy. Domyślny styl modelu jest
    proceduralny.
  \item \textbf{Format odpowiedzi to sprawa pierwszorzędna.} Przez niego lokalne modele o otwartych
    wagach okazały się dla nas bezużyteczne, a nawet model hostowany wstawił prozę do pięciu
    odpowiedzi, co spowodowało połowę wszystkich błędów kompilacji.
  \item \textbf{Zadania w stylu HumanEval są zbyt łatwe}, by porównywać paradygmaty; potrzebne są
    trudniejsze, precyzyjnie wyspecyfikowane zadania.
\end{itemize}

\section{Dalsze prace}

\begin{enumerate}
  \item \textbf{Poprawienie benchmarku.} Poprawić predykat pomocniczy budujący łańcuch w Prologu
    (zadanie 2), udokumentować nazwy i tryby argumentów w kontrakcie zadania 1 dla Prologu, sprawić,
    by program testowy w Smalltalku raportował komunikaty wyjątków, i ponownie uruchomić dotknięte
    zadania.
  \item \textbf{Powtórzone próbki.} Uruchomić od pięciu do dziesięciu próbek na zadanie, aby
    oszacować pass@$k$ i przedziały ufności, a wyniki analizować z zadaniem i językiem jako efektami
    losowymi i stałymi.
  \item \textbf{Więcej modeli.} Ocenić większe modele (Sonnet, Opus) i inne rodziny modeli, w tym
    modele o otwartych wagach udostępniane przez hostowanego dostawcę inferencji, co pozwala uniknąć
    problemów z wydajnością wdrożenia lokalnego.
  \item \textbf{Walidacja sędziego.} Zebrać oceny ludzkich ekspertów dla warstwowej próbki
    rozwiązań, użyć sędziego z innej rodziny modeli i porównać oceny bezwzględne z preferencjami
    parami.
  \item \textbf{Zadania natywne dla paradygmatów.} Dodać zadania, których naturalnym rozwiązaniem
    jest model obiektowy lub zbiór relacji (symulacje, silniki reguł, problemy z ograniczeniami),
    a także kolejne paradygmaty lub języki (np.\ OCaml, Erlang/Elixir, Rust).
  \item \textbf{Ablacje promptów.} Zmierzyć, w jakim stopniu krótkie wskazówki dotyczące paradygmatu
    w promptach (\cref{tab:guidance}) zmieniają idiomatyczność, oraz czy ściślejszy kontrakt
    odpowiedzi (np.\ pole z kodem w ustrukturyzowanej odpowiedzi) eliminuje błędy formatu.
\end{enumerate}
